\documentclass[aspectratio=169,11pt]{beamer}

\usepackage[utf8]{inputenc}
\usepackage[T1]{fontenc}
\usepackage[brazil]{babel}
\usepackage{lmodern}
\usepackage{tikz}
\usetikzlibrary{arrows.meta,positioning,shapes.geometric}
\usepackage{hyperref}

\usetheme{Madrid}
\setbeamertemplate{navigation symbols}{}
\setbeamertemplate{footline}[frame number]

% Paleta inspirada na identidade visual FIAP: vermelho, preto e branco.
\definecolor{fiapred}{HTML}{ED1C24}
\definecolor{fiapblack}{HTML}{171717}
\definecolor{fiapgray}{HTML}{292929}
\definecolor{fiaplight}{HTML}{353535}
\definecolor{fiapwhite}{HTML}{FFFFFF}

\setbeamercolor{normal text}{fg=fiapwhite,bg=fiapblack}
\setbeamercolor{structure}{fg=fiapred}
\setbeamercolor{palette primary}{fg=fiapwhite,bg=fiapgray}
\setbeamercolor{palette secondary}{fg=fiapwhite,bg=fiapblack}
\setbeamercolor{palette tertiary}{fg=fiapwhite,bg=fiapred}
\setbeamercolor{palette quaternary}{fg=fiapwhite,bg=fiapgray}
\setbeamercolor{title}{fg=fiapwhite}
\setbeamercolor{subtitle}{fg=fiapwhite}
\setbeamercolor{author}{fg=fiapwhite}
\setbeamercolor{date}{fg=fiapwhite}
\setbeamercolor{frametitle}{fg=fiapwhite,bg=fiapgray}
\setbeamercolor{block title}{fg=fiapwhite,bg=fiapred}
\setbeamercolor{block body}{fg=fiapwhite,bg=fiaplight}
\setbeamercolor{alertblock title}{fg=fiapwhite,bg=fiapred}
\setbeamercolor{alertblock body}{fg=fiapwhite,bg=fiaplight}
\setbeamercolor{footline}{fg=fiapwhite,bg=fiapblack}
\setbeamercolor{item}{fg=fiapred}
\setbeamercolor{itemize item}{fg=fiapred}
\setbeamercolor{itemize subitem}{fg=fiapred}
\setbeamercolor{description item}{fg=fiapwhite}
\hypersetup{colorlinks=true,linkcolor=fiapred,urlcolor=fiapred}

\tikzset{
  flowbox/.style={rectangle, rounded corners=2pt, draw=fiapred,
    fill=fiaplight, text=fiapwhite, thick, align=center, minimum height=0.9cm,
    text width=2.35cm, font=\small},
  audioflow/.style={rectangle, rounded corners=2pt, draw=fiapblack,
    fill=fiaplight, text=fiapwhite, thick, align=center, minimum height=0.9cm,
    text width=2.35cm, font=\small},
  flowarrow/.style={-{Stealth[length=2.5mm]}, thick, draw=fiapred}
}

\title{Pipelines de análise de vídeo e áudio}
\subtitle{Detecção de instrumentos e análise de fala}
\author{Tech Challenge — Fase 4}
\date{Novembro de 2026}

\begin{document}

% Slide 1 — abertura
\begin{frame}
  \titlepage
\end{frame}

% Slide 2 — visão dos dois pipelines
\begin{frame}{Dois pipelines independentes}
  \begin{columns}[T,onlytextwidth]
    \column{0.48\textwidth}
      \begin{block}{Visão computacional — \texttt{video\_review}}
        Rastreamento e detecção de instrumentos em vídeo cirúrgico com OpenCV e YOLOv8. Gera previsões e artefatos para revisão humana.
      \end{block}
    \column{0.48\textwidth}
      \begin{block}{Áudio — \texttt{therapy\_audio}}
        Extrai e transcreve áudio em português; classifica a polaridade textual com AWS Comprehend e permite análise complementar pela OpenAI.
      \end{block}
  \end{columns}
  \vspace{0.5em}
  \begin{alertblock}{Escopo}
    Protótipos acadêmicos. Não realizam diagnóstico nem substituem revisão humana ou avaliação independente.
  \end{alertblock}
\end{frame}

% Slide 3 — Cholec80
\begin{frame}{Dataset do pipeline de detecção: Cholec80}
  \begin{itemize}
    \item 80 vídeos de colecistectomias laparoscópicas realizadas por 13 cirurgiões.
    \item Vídeos a 25 fps; fases cirúrgicas anotadas por quadro e presença binária de instrumentos a 1 fps.
    \item Rótulos de presença não fornecem, por si só, caixas delimitadoras. As caixas usadas no projeto são rastreadas/revisadas e precisam de validação.
    \item Licença informada pela CAMMA: CC BY-NC-SA 4.0; conferir as condições de acesso, uso e citação antes de reutilizar ou redistribuir.
  \end{itemize}
  \vfill
  {\scriptsize
  Fonte e acesso: \url{https://camma.unistra.fr/datasets/}\\
  Twinanda et al., \textit{EndoNet}, IEEE TMI, 2017: \url{https://doi.org/10.1109/TMI.2016.2593957}}
\end{frame}

% Slide 4 — fluxo de detecção
\begin{frame}{Fluxo do pipeline de detecção de vídeo}
  \centering
  \begin{tikzpicture}[node distance=0.65cm and 0.42cm]
    \node[flowbox] (video) {Vídeo cirúrgico\\Cholec80};
    \node[flowbox, right=of video] (track) {Selecionar ferramenta\\e rastrear (CSRT/KCF)};
    \node[flowbox, right=of track] (review) {Revisar caixas\\e validar rótulos};
    \node[flowbox, right=of review] (dataset) {Preparar dataset\\YOLO};
    \node[flowbox, below=1.05cm of dataset] (train) {Treinar YOLOv8\\e acompanhar validação};
    \node[flowbox, left=of train] (infer) {Inferir em vídeo\\e exportar CSV};
    \node[flowbox, left=of infer] (report) {Relatório da primeira\\detecção estável};
    \node[flowbox, left=of report] (human) {Revisão\\humana};
    \draw[flowarrow] (video) -- (track);
    \draw[flowarrow] (track) -- (review);
    \draw[flowarrow] (review) -- (dataset);
    \draw[flowarrow] (dataset) -- (train);
    \draw[flowarrow] (train) -- (infer);
    \draw[flowarrow] (infer) -- (report);
    \draw[flowarrow] (report) -- (human);
  \end{tikzpicture}
  \vfill
  {\scriptsize O split atual é por imagem; não comprova generalização para vídeos/procedimentos independentes.}
\end{frame}

% Slide 5 — fonte do vídeo de terapia
\begin{frame}{Vídeo usado no pipeline de áudio}
  \begin{block}{Demonstração publicada}
    \textit{Paulo Knapp Demonstra uma Sessão de Terapia Cognitivo-Comportamental — TCC}, publicado pelo canal Centro de Mindfulness. A descrição informa que foi filmado no Congresso CCE de 2010.
  \end{block}
  \begin{itemize}
    \item O projeto processa um arquivo local; o caminho de entrada é configurável.
    \item A publicação é identificada como demonstração. Não há base para descrevê-la como gravação de atendimento clínico real.
    \item Antes de baixar, processar ou redistribuir, verificar licença, autorização e privacidade.
  \end{itemize}
  \vfill
  {\scriptsize Requisito do projeto e vídeo: \url{https://www.youtube.com/watch?v=VS6YMAjdGZA}\\
  Origem documentada em: \texttt{documentos/requisitos/pipeline-audio-sessao-terapia.md}}
\end{frame}

% Slide 6 — fluxo de áudio
\begin{frame}{Fluxo do pipeline de áudio}
  \centering
  \begin{tikzpicture}[node distance=0.72cm and 0.5cm]
    \node[audioflow] (input) {Vídeo local\\e intervalo};
    \node[audioflow, right=of input] (extract) {Extrair áudio\\WAV mono/16 kHz};
    \node[audioflow, right=of extract] (segment) {Segmentar\\próximo a pausas};
    \node[audioflow, right=of segment] (speech) {Google Speech\\transcrição pt-BR};
    \node[audioflow, below=1.15cm of speech] (aws) {AWS Comprehend\\polaridade por trecho};
    \node[audioflow, left=of aws] (openai) {OpenAI\\opcional};
    \node[audioflow, left=of openai] (output) {Exibir resultados\\e salvar WAV/texto};
    \node[audioflow, left=of output] (review) {Revisão\\humana};
    \draw[flowarrow] (input) -- (extract);
    \draw[flowarrow] (extract) -- (segment);
    \draw[flowarrow] (segment) -- (speech);
    \draw[flowarrow] (speech) -- (aws);
    \draw[flowarrow] (aws) -- (openai);
    \draw[flowarrow] (openai) -- (output);
    \draw[flowarrow] (output) -- (review);
  \end{tikzpicture}
  \vfill
  {\scriptsize Polaridade é classificação textual, não avaliação clínica ou medida de intensidade emocional.}
\end{frame}

% Slide 7 — demonstração
\begin{frame}[plain]
  \centering
  \vfill
  {\Huge Demonstração ao vivo}\par
  \vspace{1em}
  {\large Dois fluxos independentes; resultados para revisão humana.}\par
  \vfill
\end{frame}

\end{document}
